\renewcommand{\chaptertitlename}{Chapter}
\chapter{Star: Are you okey?}
\renewcommand{\chaptertitlename}{Capítulo}

\nowplaying{KILLSWITCH (Levia Remix)}{mekaloton \& Levia}

\lettrine{B}{usco} tu mirada entre el polvo y el ruido. —¿Estás bien? Hey, ¿qué tal, mi pequeño felino? ¿Te encuentras bien?

El cielo se parte en pedazos de fuego: meteoritos con forma de estrella caen hacia nosotros a una velocidad que no debería existir. El rostro de León cambia, y por primera vez desde que lo conozco, veo miedo dibujado en él.

\illustration{images/Chapter-1/leon-asustado.png}{León, aterrado, ve caer los meteoritos por primera vez.}

—¿Qué sucede? ¿Por qué tu rostro cambió así? ¿Por qué tienes miedo?

No hay tiempo para respuesta. Uno de los meteoritos cae directo hacia mí y León me empuja con todo su peso, apartándome del impacto.

\illustration{images/Chapter-1/leon-super-asustado.png}{El instante justo después de empujar a Jhonny fuera de peligro.}

—¡Auch! ¿Por qué hiciste eso? —Ya entiendo. Esto no va a parar. Siguen cayendo, una tras otra. Rápido, corre conmigo.

Tal vez, para ayudar a otros, deba ayudarme primero a mí mismo.

A lo lejos, entre gritos, alguien menciona un «amuleto encantado». Esquivamos como podemos, con una coreografía improvisada de saltos y giros que se siente casi elegante, casi \emph{cool}, aunque yo todavía no tengo ningún poder. Solo sé esquivar.

Wow. Así que esta es mi condición física real. ¿Por qué no me ejercité cuando pude? ¿Por qué no cuidé mi cuerpo cuando era el momento? ¿Qué edad tengo? ¿En dónde estoy?

Alguien pasa corriendo a una velocidad imposible y me grita sin detenerse: \emph{``You're too slow.''}

Una chica corre a mi lado, hablando más para sí misma que para mí: —Esto no hubiera pasado si hubiera continuado mis estudios en la San Marino.

Tropiezo. León me carga sobre su espalda sin detener el paso. Algo me lanza una telaraña directo a la boca; la escupo, desorientado. ¿Dónde estoy?

Corremos en dirección contraria a los meteoritos, hacia un muro enorme que divide lo que aún se puede cruzar de lo que ya no.

A mi izquierda, un niño le grita a su robot que busque en su bolsillo algo que los pueda sacar con vida de ahí. El robot obedece: de un bolsillo diminuto empiezan a salir objetos que no deberían caber, uno tras otro, imposibles.

A mi derecha, un ser dibuja círculos en el suelo junto a una armadura metálica que convierte el piso mismo en un vehículo.

Todo el mundo está improvisando su propia forma de sobrevivir. Yo solo sé correr. Y esquivar. Y preguntarme, mientras el cielo sigue cayendo, cuánto tiempo más podré seguir haciendo solamente eso.
